\sgaStableTarget{sga6.e08.prop.5-2}\textbf{Proposition 5.2.} \uline{We have on} @@M0@@ \uline{the exact sequence}:
@@M1@@
\uline{i}.\uline{e}. \uline{we have an isomorphism}
@@M2@@

Consider the vector bundle over @@M3@@ defined by @@M4@@, and denote by @@M5@@ the complement of the zero section of this bundle. We have the commutative diagram:
@@M6@@
The morphism @@M7@@ identifies @@M8@@ with the affine @@M9@@-scheme defined by the sheaf of algebras @@M10@@; @@M11@@ is a graded @@M12@@-Algebra.

Since @@M13@@ are smooth morphisms, one may write the exact sequence of sheaves of relative differentials; since @@M14@@ is affine over @@M15@@, one may write this sequence in terms of @@M16@@-Modules; one thus obtains:
\sgaStableTarget{sga6.e08.eq.5-2-1}@@M17@@

Now this exact sequence is a graded exact sequence of graded @@M18@@-Modules. Indeed, it follows from the commutative diagram:
@@M19@@
where the arrows are algebra homomorphisms compatible with the grading of @@M20@@ (grading defined by degree in the indeterminate @@M21@@), the two other algebras being supposed concentrated
